\section{Erdős \#389}

\textbf{Problem.} Is it true that for every $n \geq 1$ there exists $k \geq 1$ such that
$n(n+1)\cdots(n+k-1) \mid (n+k)(n+k+1)\cdots(n+2k-1)$?

\textbf{What was attempted.} For each $n$ from $1$ to $200$, a brute-force search was run to find the smallest $k$ (up to a cutoff) for which the divisibility holds. The ratio was identified as
\[
\frac{\prod_{i=0}^{k-1}(n+k+i)}{\prod_{i=0}^{k-1}(n+i)} = \frac{\binom{n+2k-1}{k}}{\binom{n+k-1}{k}},
\]
and the case $k = n$ was tested as a candidate witness.

\textbf{What the computation showed.} For $n = 1$ through $5$, valid $k$ values were found (e.g., $n=4, k=207$; $n=5, k=206$), but these required a cutoff far exceeding the initial bound of $300$. The candidate $k = n$ fails for all $n \geq 2$. For $n = 6$, no valid $k$ was found up to $k = 2000$; for $n = 7$ a solution $k = 984$ appeared, but $n = 8, 9, 10$ again failed within the same range. The search also aborted early due to a Python syntax error (\texttt{2k} instead of \texttt{2*k}) in round~1, corrupting the initial sweep entirely.

\textbf{What went wrong.} The computation is inconclusive on multiple grounds: (i) the initial syntax error invalidated the first pass; (ii) the search cutoff was far too small---solutions may exist for $n = 6, 8, 9, 10$ at $k \gg 2000$, as the trend for small $n$ already shows $k$ growing erratically into the hundreds; (iii) no asymptotic or number-theoretic analysis was completed to explain why solutions are sparse or to bound the required $k$. The observation that no $k \leq 2000$ works for $n = 6$ is not evidence against the conjecture; it is merely evidence that the required $k$ is large.

\textbf{Verdict:} No counterexample found within the (severely limited) search range, consistent with the conjecture but constituting no progress on it.